\begin{beginnerstatement}[Kurz erklärt: Template-Strategien]

Ein Trade-off ist eine Abwägung, bei der ein Vorteil mit einem Nachteil
verbunden ist.

Ein Framework-Starter bietet einen einfachen Einstieg für einen eng
begrenzten Technikbereich, etwa nur für das Frontend.

Ein monolithisches Template enthält alle vorgesehenen Bausteine gemeinsam,
auch wenn ein kleines Projekt einige davon nicht benötigt.

Bei einem Copy-Paste-Starter wird die Ausgangsbasis einmal kopiert; spätere
Verbesserungen der Vorlage gelangen nicht automatisch in das Projekt.

Die Profilstrategie der Fallstudie wählt Bausteine vor der Generierung aus,
nimmt dafür aber zusätzliche Pflege- und Auswahlkomplexität in Kauf.

\end{beginnerstatement}
